% Chapter 6

\chapter{Limitations and Future Work}\label{Chapter6}

%----------------------------------------------------------------------------------------
%	SECTION 1
%----------------------------------------------------------------------------------------

\section{Limitations}\label{sec:limits}

Several limitations constrain the interpretation and generalisability of these results.

The most immediate limitation is the single-centre design. All data originate from one rehabilitation facility, so selection bias is possible and generalisation to other clinical populations cannot be assumed \parencite{corrigan2010epidemiology}. The use of established neuropsychological tests, including Trail Making, digit span, Rey Auditory Verbal Learning, orientation, and language measures, partly mitigates this concern because these instruments have long clinical histories and international validation \parencite{lezak2012neuropsychological}. The six cognitive domains also provide a standard framework for describing post-injury function. Multi-centre validation is nevertheless required before the severity structure can be interpreted as a general feature of ABI rather than a property of this service.

The structure recovered here is also a limitation on interpretation. The main finding is a continuum, not a set of natural categories. The cohort falls along a single overall-severity dimension, with PC1 explaining 56\% of the variance, rather than into clearly dissociable profiles. The three tiers are useful as a discretisation of that gradient, but they should not be treated as phenotypes with hard boundaries. Tier count depends partly on clustering resolution: the primary MICE pipeline produces three, other imputation methods produce three to five, and the result shifts with the missingness threshold. Assessments near boundaries require caution. A continuous severity score may ultimately be more faithful than discrete bands (Section~\ref{sec:future}).

The analysis is also limited in its ability to model longitudinal change. ABI recovery is often fastest in the first months and then continues more slowly over years \parencite{ponsford2014longitudinal, maas2017traumatic}. A single assessment captures only one point in that process. The dataset does include multiple assessments for many patients, which supports the within-patient analysis reported below, but the timing of follow-up is still clinically driven rather than standardised.

Among the 5{,}206 patients with longitudinal data (Section~\ref{sec:robustness}), the pattern is clinically plausible: 50.2\% remain in the same tier, and among those who change tier, 706 move toward better function compared with 206 who move toward worse function. This 3:1 ratio matches the expected direction of recovery. However, it cannot fully separate genuine recovery from selection effects, such as patients being discharged once they improve. A future longitudinal design with harmonised assessment intervals is needed for that question.

The imputation strategy imposes additional constraints. The missing-at-random (MAR) assumption cannot be tested from observed data alone \parencite{rubin1976inference, little2019statistical}. In clinical settings, a test may be incomplete precisely because the patient cannot finish it, so the missing value may be informative and may lie in the impaired range. If MNAR patterns of this kind were dominant, function could be overestimated. The conservative tipping-point check (Section~\ref{sec:robustness}) re-imputes every missing value to the worst observed performance. Under that extreme assumption, the Global Impairment tier persists and retains 65\% of its members, although the finer partition degrades. A formal delta-shifted sensitivity analysis would refine this bound. The structured co-occurrence of missingness is somewhat reassuring because it points to test-battery protocols rather than purely outcome-based dropout. The comparison of 10 imputation methods also functions as a sensitivity analysis: 86.2\% of assessments receive high-confidence consensus labels, suggesting that plausible imputation differences do not alter the main severity structure for most of the cohort.

A related caveat is the Tier-1 filter, which excludes assessments with more than half their values missing. On inspection (Section~\ref{sec:robustness}), excluded records completed only 0.8 of 14 tests on average and were demographically very similar to retained records. The filter therefore appears to remove near-empty administrative records rather than a hidden severe subgroup. Still, the cognitive severity of excluded records cannot be observed directly, and the 30--70\% threshold sensitivity analysis only bounds the effect of this choice. Finally, the pipeline imputes and then clusters, rather than using a model-based alternative such as latent profile analysis with full-information maximum likelihood. The Gaussian-mixture baseline (Section~\ref{sec:robustness}) shows the Bayesian Information Criterion falling monotonically with the number of classes, supporting a continuum rather than a discrete optimum. A fuller FIML-based latent-profile analysis remains a worthwhile future extension.

Feature engineering remains consequential. Non-cognitive variables are excluded, and every variable is winsorised, direction-aligned, and standardised before aggregation so that no raw scale can dominate a composite and generate an artificial profile difference. The broader methodological lesson is direct: unsupervised phenotyping depends on preprocessing. Any striking phenotype should be tested against scaling, variable-selection, and cohort-definition choices before it is accepted.

HDBSCAN also introduces hyperparameter sensitivity. I conducted a broad search ($9 \times 120$ configurations) and selected the optimal hyperparameters, including \texttt{min\_cluster\_size}$= 1000$ and \texttt{epsilon}$= 0.0$, by maximising the composite quality criterion $Q$. Because the underlying structure is a gradient, however, the number of tiers is partly tied to \texttt{min\_cluster\_size}. More broadly, UMAP+HDBSCAN is stochastic and can be sensitive to random initialisation and imperfect global-structure preservation. I mitigated this through ten random seeds per configuration, per-method bootstrap analysis, 80\% subsampling, and comparison with PCA. The headline conclusion does not depend on UMAP alone: PCA also supports the severity gradient (PC1 = 56\%). Re-running the pipeline under ten random seeds produced near-identical labels (mean pairwise ARI 0.957; Section~\ref{sec:robustness}). The exact cluster count can fluctuate, but the severity ordering is stable across these checks.

Computational cost is the final practical limitation. Running UMAP+HDBSCAN 50 times for each method is non-trivial, although the modular architecture permits parallelisation. Because the 14 variables are reduced to 6 domain scores, runtimes remained manageable on a consumer-grade laptop.

%----------------------------------------------------------------------------------------
%	SECTION 2
%----------------------------------------------------------------------------------------

\section{Future Directions}\label{sec:future}

The findings and limitations suggest four concrete directions for future work.

\textbf{Model the continuum directly.} The clearest implication of the severity-gradient finding is that a continuous cognitive-severity \emph{score}---for example, the first principal component, or a supervised index validated against functional outcome---may be a more faithful and more useful summary than a discrete tier label. Future work should compare a continuous severity index against the tiered representation for prognosis and for matching rehabilitation intensity to need, treating the tiers as an interpretable discretisation of the underlying score rather than as natural kinds.

\textbf{Longitudinal clustering} is the most direct extension. Applying the pipeline to assessments collected at fixed time points would let the analysis describe how a patient's position on the severity gradient changes over the recovery trajectory: whether the initial level predicts long-term outcome, and whether recovery rates differ across the severity range \parencite{ponsford2014longitudinal}. Trajectory-based severity modelling would be a far more useful input to prognosis and rehabilitation-intensity decisions than the cross-sectional snapshot.

\textbf{Multi-centre validation} is the prerequisite for treating the severity stratification as a feature of ABI rather than of this particular service. A multi-centre study across different healthcare systems would test whether the one-dimensional severity structure survives changes in catchment, test battery, and assessment timing, and would let researchers develop harmonisation strategies for cross-centre comparison \parencite{maas2017traumatic}.

\textbf{Neuroimaging integration} would bridge the gap between cognitive severity and neural mechanism. Linking a patient's position on the severity gradient to structural and functional brain-imaging measures could identify the neural correlates of global cognitive impairment after ABI, the kind of evidence that could guide targeted intervention \parencite{saatman2008classification, maas2017traumatic}.
